% Сообщение Ярославу: в работе уверенно выполнены задания на простую вероятность, независимые события, математическое ожидание и вычисление выражения. Ошибки и пропуски связаны прежде всего с геометрией, условной вероятностью, прикладными задачами, чтением графиков, рациональными уравнениями и неравенствами, а также заданиями повышенной сложности. В задании 6 неверно определено число участников в запасной аудитории: нужно вычесть обе основные аудитории, а не одну.
\documentclass[12pt,a4paper]{article}
\usepackage[a4paper,margin=18mm]{geometry}
\usepackage[T2A]{fontenc}
\usepackage[utf8]{inputenc}
\usepackage[russian]{babel}
\usepackage{amsmath,amssymb,mathtools}
\usepackage{array,booktabs,longtable}
\usepackage{xcolor}
\usepackage{hyperref}
\usepackage{tikz}
\hypersetup{colorlinks=true,linkcolor=blue!55!black}
\setlength{\parindent}{0pt}
\setlength{\parskip}{0.55em}
\renewcommand{\arraystretch}{1.18}
\newcommand{\errorlabel}[1]{\textcolor{red}{\textbf{Ошибка:}} #1}
\newcommand{\keyidea}[1]{\textcolor{blue!60!black}{\textbf{Ключевая идея:}} #1}
\newcommand{\finalnote}[1]{\textcolor{teal!60!black}{\textbf{Итог:}} #1}
\newcommand{\taskhead}[1]{\newpage\phantomsection\label{task:#1}\section*{Задание #1}}
\newcommand{\subhead}[1]{\vspace{0.2em}\textbf{#1}\par}
\begin{document}

\hypersetup{pageanchor=false}
\begin{titlepage}
\centering
\vspace*{0.25\textheight}
{\LARGE\bfseries Ревью домашней работы\par}
\vspace{2.2cm}
{\large Автор: Лёвин Артём Александрович\par}
\vspace{0.8cm}
{\large 22 сентября 2026 года\par}
\vfill
\begin{tikzpicture}[x=1cm,y=1cm]
\draw[blue!35!black,line width=0.7pt] (-4,0) .. controls (-2,0.7) and (2,-0.7) .. (4,0);
\end{tikzpicture}
\end{titlepage}
\hypersetup{pageanchor=true}

\newpage
\section*{Сводная таблица}
\small
\setlength{\tabcolsep}{3.5pt}
\begin{longtable}{>{\centering\arraybackslash}p{0.07\linewidth} p{0.16\linewidth} p{0.29\linewidth} p{0.18\linewidth} p{0.18\linewidth}}
\toprule
\textbf{№} & \textbf{Тема} & \textbf{Суть ошибки} & \textbf{Ваш ответ} & \textbf{Правильный ответ}\\
\midrule
\endfirsthead
\toprule
\textbf{№} & \textbf{Тема} & \textbf{Суть ошибки} & \textbf{Ваш ответ} & \textbf{Правильный ответ}\\
\midrule
\endhead
\hyperref[task:1]{1} & Окружность и касательные & Решение не записано. & не указан & $64^\circ$\\
\hyperref[task:2]{2} & Описанная окружность & Решение не записано. & не указан & $150^\circ$\\
\hyperref[task:3]{3} & Скалярное произведение & Решение не записано. & не указан & $45^\circ$\\
\hyperref[task:4]{4} & Векторы & Решение не записано. & не указан & $13$\\
\hyperref[task:6]{6} & Вероятность & Вычтена только одна основная аудитория вместо двух. & $0{,}6$ & $0{,}2$\\
\hyperref[task:8]{8} & Условная вероятность & Решение не записано. & не указан & $0{,}43$\\
\hyperref[task:10]{10} & Линейные выражения & Решение не записано. & не указан & $10$\\
\hyperref[task:12]{12} & Неравенство по формуле & Решение не записано. & не указан & $400$ К\\
\hyperref[task:13]{13} & Квадратичная функция & Решение не записано. & не указан & $110$ м\\
\hyperref[task:14]{14} & Движение & Решение не записано. & не указан & $600$ м\\
\hyperref[task:15]{15} & Движение по реке & Решение не записано. & не указан & $288$ км\\
\hyperref[task:16]{16} & График $k/x+a$ & Решение не записано. & не указан & $-\frac{74}{25}$\\
\hyperref[task:17]{17} & Квадратичная функция & Решение не записано. & не указан & $21$\\
\hyperref[task:18]{18} & Рациональное уравнение & Решение не записано. & не указан & $-1,5,7\pm2\sqrt{11}$; на $[-2;3]$: $-1,7-2\sqrt{11}$\\
\hyperref[task:19]{19} & Рациональное неравенство & Решение не записано. & не указан & $(-\infty;-4]\cup(-3;-1]\cup(0;7)\cup(7;+\infty)$\\
\hyperref[task:20]{20} & Геометрия окружностей & Решение не записано. & не указан & $AE\parallel BD$, $AC=\frac{240}{17}$\\
\hyperref[task:21]{21} & Параметр & Решение не записано. & не указан & $3<a<\frac{13}{4}$\\
\hyperref[task:22]{22} & Делимость & Решение не записано. & не указан & а) нет; б) нет; в) $21$\\
\bottomrule
\end{longtable}
\setlength{\tabcolsep}{6pt}
\normalsize

\newpage
\thispagestyle{empty}
\vspace*{\fill}
\begin{center}{\LARGE\bfseries Разбор ошибок}\end{center}
\vspace*{\fill}

\taskhead{1}
\subhead{Текст задания}
Через концы $A$ и $B$ дуги окружности с центром $O$ проведены касательные $AC$ и $BC$. Угол $CAB=32^\circ$. Найдите $\angle AOB$.
\subhead{Ваше решение} Прочерк. \subhead{Ваш ответ} Не указан.
\errorlabel{решение отсутствует.}
\subhead{Где именно возникает ошибка}
Решение не начато. Не использованы равенство касательных $CA=CB$ и перпендикулярность радиусов касательным.
\subhead{Почему это неверно}
Из $CA=CB$ следует $\angle CAB=\angle CBA=32^\circ$, поэтому $\angle ACB=116^\circ$. В четырёхугольнике $AOCB$ углы при $A$ и $B$ прямые, значит $\angle AOB+\angle ACB=180^\circ$.
\keyidea{сначала найти угол между касательными, затем центральный угол.}
\subhead{Исправление и полное решение}
\[
\angle AOB=180^\circ-(180^\circ-2\cdot32^\circ)=64^\circ.
\]
\subhead{Проверка}
$116^\circ+64^\circ=180^\circ$, что соответствует теореме об угле между касательными.
\finalnote{ответ: $64^\circ$.}
\subhead{Задача на закрепление}
Через точки $A$ и $B$ окружности проведены касательные, пересекающиеся в $C$. Если $\angle CAB=27^\circ$, найдите $\angle AOB$.

\taskhead{2}
\subhead{Текст задания}
Сторона $AB$ треугольника $ABC$ с тупым углом $C$ равна радиусу описанной окружности. Найдите $C$.
\subhead{Ваше решение} Прочерк. \subhead{Ваш ответ} Не указан.
\errorlabel{решение отсутствует.}
\subhead{Где именно возникает ошибка}
Не применена расширенная теорема синусов.
\subhead{Почему это неверно}
Сторона $AB$ лежит напротив угла $C$, поэтому $AB=2R\sin C$. По условию $AB=R$.
\keyidea{получить значение синуса и выбрать именно тупой угол.}
\subhead{Исправление и полное решение}
\[
R=2R\sin C,\qquad \sin C=\frac12.
\]
В промежутке $(0^\circ;180^\circ)$ это даёт $30^\circ$ или $150^\circ$. Угол тупой, значит
\[
C=150^\circ.
\]
\subhead{Проверка}
$2R\sin150^\circ=R$.
\finalnote{ответ: $150^\circ$.}
\subhead{Задача на закрепление}
В треугольнике $ABC$ сторона $AB=\sqrt3\,R$, где $R$ --- радиус описанной окружности. Угол $C$ тупой. Найдите $C$.

\taskhead{3}
\subhead{Текст задания}
$\vec a=(1;1)$, $\vec b=(0;2)$. Используя скалярное произведение, найдите угол между векторами.
\subhead{Ваше решение} Прочерк. \subhead{Ваш ответ} Не указан.
\errorlabel{решение отсутствует.}
\subhead{Где именно возникает ошибка}
Не вычислены скалярное произведение и длины векторов.
\subhead{Почему это неверно}
Используется формула $\vec a\cdot\vec b=|\vec a||\vec b|\cos\varphi$.
\keyidea{найти косинус угла через координаты.}
\subhead{Исправление и полное решение}
\[
\vec a\cdot\vec b=2,\quad |\vec a|=\sqrt2,\quad |\vec b|=2,
\]
\[
\cos\varphi=\frac{2}{2\sqrt2}=\frac{\sqrt2}{2},
\qquad \varphi=45^\circ.
\]
\subhead{Проверка}
Вектор $(1;1)$ образует с осью $x$ угол $45^\circ$, а $(0;2)$ направлен вертикально.
\finalnote{ответ: $45^\circ$.}
\subhead{Задача на закрепление}
Даны $\vec a=(2;2)$ и $\vec b=(0;3)$. Используя скалярное произведение, найдите угол между ними.

\taskhead{4}
\subhead{Текст задания}
$\vec a=(1;1)$, $\vec b=(0;7)$. Найдите длину $5\vec a+\vec b$.
\subhead{Ваше решение} Прочерк. \subhead{Ваш ответ} Не указан.
\errorlabel{решение отсутствует.}
\subhead{Где именно возникает ошибка}
Не выполнено координатное сложение.
\subhead{Почему это неверно}
Сначала нужно получить координаты итогового вектора, затем найти его длину.
\keyidea{координаты, затем теорема Пифагора.}
\subhead{Исправление и полное решение}
\[
5\vec a+\vec b=(5;5)+(0;7)=(5;12),
\]
\[
|5\vec a+\vec b|=\sqrt{5^2+12^2}=13.
\]
\subhead{Проверка}
$13^2=5^2+12^2$.
\finalnote{ответ: $13$.}
\subhead{Задача на закрепление}
Даны $\vec a=(2;-1)$ и $\vec b=(1;4)$. Найдите длину $3\vec a+\vec b$.

\taskhead{6}
\subhead{Текст задания}
Из $450$ участников в первых двух аудиториях разместили по $180$ человек, остальных --- в запасной. Найдите вероятность попасть в запасную аудиторию.
\subhead{Ваше решение}
$450-180=270$, затем $\frac{270}{450}=0{,}6$.
\subhead{Ваш ответ} $0{,}6$.
\errorlabel{вычтена только одна из двух основных аудиторий.}
\subhead{Где именно возникает ошибка}
Первый неверный шаг --- $450-180=270$ как число участников запасной аудитории. В двух основных аудиториях вместе $180+180=360$ человек.
\subhead{Почему это неверно}
Благоприятных исходов не $270$, а $450-360=90$.
\keyidea{сначала посчитать оставшихся, затем разделить на общее число участников.}
\subhead{Исправление и полное решение}
\[
450-180-180=90,\qquad
P=\frac{90}{450}=\frac15=0{,}2.
\]
\subhead{Проверка}
$90$ человек --- пятая часть от $450$.
\finalnote{ответ: $0{,}2$.}
\subhead{Задача на закрепление}
На конференции $520$ участников. В двух основных залах разместили по $200$ человек, остальных --- в запасном. Найдите вероятность попасть в запасной зал.

\taskhead{8}
\subhead{Текст задания}
Чувствительность теста $86\%$, специфичность $94\%$, общая доля положительных тестов $10\%$. При положительном тесте найдите вероятность заболевания.
\subhead{Ваше решение} Прочерк. \subhead{Ваш ответ} Не указан.
\errorlabel{решение отсутствует.}
\subhead{Где именно возникает ошибка}
Не составлено уравнение полной вероятности положительного теста.
\subhead{Почему это неверно}
Положительный результат возникает у заболевших с вероятностью $0{,}86$ и у незаболевших как ложноположительный с вероятностью $0{,}06$.
\keyidea{сначала найти долю заболевших, затем условную вероятность.}
\subhead{Исправление и полное решение}
Пусть $p$ --- доля заболевших. Тогда
\[
0{,}86p+0{,}06(1-p)=0{,}10,
\]
\[
0{,}80p=0{,}04,\qquad p=0{,}05.
\]
Совместная вероятность заболевания и положительного теста:
\[
0{,}05\cdot0{,}86=0{,}043.
\]
Поэтому
\[
P(\text{болен}\mid\text{тест }+)=\frac{0{,}043}{0{,}10}=0{,}43.
\]
\subhead{Проверка}
Из ста обследованных в среднем $4{,}3$ человека одновременно больны и имеют положительный тест, а положительных тестов всего $10$.
\finalnote{ответ: $0{,}43$.}
\subhead{Задача на закрепление}
Чувствительность теста $90\%$, специфичность $95\%$, общая доля положительных результатов $14\%$. Найдите вероятность заболевания при положительном результате.

\taskhead{10}
\subhead{Текст задания}
Найдите $61a-11b+50$, если
\[
\frac{2a-7b+5}{7a-2b+5}=9.
\]
\subhead{Ваше решение} Прочерк. \subhead{Ваш ответ} Не указан.
\errorlabel{решение отсутствует.}
\subhead{Где именно возникает ошибка}
Не преобразовано исходное отношение в линейное равенство.
\subhead{Почему это неверно}
Отдельно находить $a$ и $b$ не требуется: после раскрытия скобок сразу получается нужная комбинация.
\keyidea{перемножить крест-накрест и выделить $61a-11b$.}
\subhead{Исправление и полное решение}
\[
2a-7b+5=9(7a-2b+5),
\]
\[
2a-7b+5=63a-18b+45,
\]
\[
61a-11b+40=0.
\]
Значит, $61a-11b=-40$, поэтому
\[
61a-11b+50=10.
\]
\subhead{Проверка}
Значение нужного выражения определяется непосредственно условием.
\finalnote{ответ: $10$.}
\subhead{Задача на закрепление}
Найдите $7a-2b+6$, если $\frac{a-2b+3}{2a-b+1}=4$.

\taskhead{12}
\subhead{Текст задания}
\[
\eta=\frac{T_1-T_2}{T_1}\cdot100\%.
\]
Найдите минимальное $T_1$, при котором $\eta\ge15\%$, если $T_2=340$ К.
\subhead{Ваше решение} Прочерк. \subhead{Ваш ответ} Не указан.
\errorlabel{решение отсутствует.}
\subhead{Где именно возникает ошибка}
Не составлено неравенство по условию «не меньше».
\subhead{Почему это неверно}
Температура в Кельвинах положительна, поэтому на $T_1$ можно умножать без смены знака.
\keyidea{подставить $T_2$ и решить неравенство.}
\subhead{Исправление и полное решение}
\[
\frac{T_1-340}{T_1}\ge0{,}15,
\]
\[
T_1-340\ge0{,}15T_1,
\qquad 0{,}85T_1\ge340,
\]
\[
T_1\ge400.
\]
Минимум равен $400$ К.
\subhead{Проверка}
При $T_1=400$ К получаем ровно $15\%$.
\finalnote{ответ: $400$ К.}
\subhead{Задача на закрепление}
Если $T_2=300$ К, найдите минимальный $T_1$, при котором КПД не меньше $20\%$.

\taskhead{13}
\subhead{Текст задания}
\[
y=-\frac{x^2}{110}+\frac{13}{11}x.
\]
Стена высотой $19$ м. Найдите наибольшее расстояние, при котором камень проходит не менее чем на $1$ м выше стены.
\subhead{Ваше решение} Прочерк. \subhead{Ваш ответ} Не указан.
\errorlabel{решение отсутствует.}
\subhead{Где именно возникает ошибка}
Не переведено условие в неравенство $y\ge20$.
\subhead{Почему это неверно}
В точке стены высота камня должна быть не меньше $19+1=20$ м.
\keyidea{решить квадратное неравенство и взять правую границу.}
\subhead{Исправление и полное решение}
\[
-\frac{x^2}{110}+\frac{13}{11}x\ge20
\]
\[
-x^2+130x\ge2200
\]
\[
x^2-130x+2200\le0.
\]
Корни:
\[
D=130^2-8800=8100,\qquad
x=\frac{130\pm90}{2},
\]
то есть $20$ и $110$. Следовательно,
\[
20\le x\le110.
\]
\subhead{Проверка}
$y(110)=-110+130=20$.
\finalnote{ответ: $110$ м.}
\subhead{Задача на закрепление}
Траектория $y=-\frac{x^2}{100}+x$. Стена высотой $15$ м. Найдите наибольшее расстояние, чтобы камень прошёл не менее чем на $1$ м выше стены.

\taskhead{14}
\subhead{Текст задания}
Поезд со скоростью $60$ км/ч полностью проходит мимо лесополосы длиной $400$ м за $1$ минуту. Найдите длину поезда.
\subhead{Ваше решение} Прочерк. \subhead{Ваш ответ} Не указан.
\errorlabel{решение отсутствует.}
\subhead{Где именно возникает ошибка}
Не учтено, что для полного прохождения объекта поезд проходит сумму своей длины и длины объекта.
\subhead{Почему это неверно}
$60$ км/ч --- это $1000$ м/мин.
\keyidea{из полного пути вычесть длину лесополосы.}
\subhead{Исправление и полное решение}
За минуту поезд проходит $1000$ м. Если его длина $L$, то
\[
L+400=1000,
\qquad L=600.
\]
\subhead{Проверка}
$600+400=1000$ м.
\finalnote{ответ: $600$ м.}
\subhead{Задача на закрепление}
Поезд движется со скоростью $72$ км/ч и полностью проходит тоннель длиной $500$ м за $45$ секунд. Найдите длину поезда.

\taskhead{15}
\subhead{Текст задания}
Собственная скорость теплохода $15$ км/ч, течение $3$ км/ч, стоянка $5$ часов, весь рейс туда и обратно $25$ часов. Найдите общий путь.
\subhead{Ваше решение} Прочерк. \subhead{Ваш ответ} Не указан.
\errorlabel{решение отсутствует.}
\subhead{Где именно возникает ошибка}
Не выделено время движения и не учтены разные скорости по и против течения.
\subhead{Почему это неверно}
Скорости равны $18$ и $12$ км/ч, а движение занимает $20$ часов.
\keyidea{составить уравнение по времени для одного и того же расстояния $S$.}
\subhead{Исправление и полное решение}
\[
\frac{S}{18}+\frac{S}{12}=20,
\]
\[
\frac{5S}{36}=20,\qquad S=144.
\]
Общий путь:
\[
2S=288\text{ км}.
\]
\subhead{Проверка}
$144/18=8$ ч, $144/12=12$ ч; $8+12+5=25$ ч.
\finalnote{ответ: $288$ км.}
\subhead{Задача на закрепление}
Собственная скорость теплохода $18$ км/ч, течение $2$ км/ч, стоянка $4$ часа, весь рейс $13$ часов. Найдите общий путь.

\taskhead{16}
\subhead{Текст задания}
По графику $f(x)=\frac{k}{x}+a$ найдите $f(50)$.
\subhead{Ваше решение}
Решение по графику не записано.
\subhead{Ваш ответ} Не указан.
\errorlabel{не найдены параметры функции по графику.}
\subhead{Где именно возникает ошибка}
Не использованы горизонтальная асимптота $y=-3$ и отмеченная точка $(-2;-4)$.
\subhead{Почему это неверно}
Для функции $\frac{k}{x}+a$ горизонтальная асимптота равна $y=a$.
\keyidea{найти $a$ по асимптоте и $k$ по точке графика.}
\subhead{Исправление и полное решение}
\[
a=-3.
\]
Точка $(-2;-4)$ даёт
\[
-4=\frac{k}{-2}-3,
\qquad k=2.
\]
Тогда
\[
f(50)=\frac{2}{50}-3=-\frac{74}{25}.
\]
\subhead{Проверка}
$-\frac{74}{25}=-2{,}96$, то есть значение немного выше асимптоты $-3$, как и должно быть при большом положительном $x$.
\finalnote{ответ: $-\frac{74}{25}$.}
\subhead{Задача на закрепление}
$f(x)=\frac{k}{x}+a$ имеет горизонтальную асимптоту $y=2$ и проходит через $(3;4)$. Найдите $f(12)$.

\taskhead{17}
\subhead{Текст задания}
По графику $f(x)=ax^2+bx+c$, где $a,b,c$ --- целые, найдите $f(2)$.
\subhead{Ваше решение} Решение не записано. \subhead{Ваш ответ} Не указан.
\errorlabel{решение отсутствует.}
\subhead{Где именно возникает ошибка}
Не использованы нули графика $-5$ и $-1$ и значение $f(0)=5$.
\subhead{Почему это неверно}
По двум корням функция имеет вид $a(x+5)(x+1)$, а третья точка задаёт $a$.
\keyidea{восстановить квадратный трёхчлен по графику.}
\subhead{Исправление и полное решение}
\[
f(x)=a(x+5)(x+1).
\]
Так как $f(0)=5$,
\[
5=5a,\qquad a=1.
\]
Следовательно,
\[
f(2)=7\cdot3=21.
\]
\subhead{Проверка}
$f(-3)=(-3+5)(-3+1)=-4$, что совпадает с вершиной на графике.
\finalnote{ответ: $21$.}
\subhead{Задача на закрепление}
Квадратичная функция имеет корни $-4$ и $2$ и проходит через $(0;-8)$. Найдите $f(3)$.

\taskhead{18}
\subhead{Текст задания}
Решите
\[
\frac{(x-1)^2}{8}+\frac8{(x-1)^2}
=7\left(\frac{x-1}{4}-\frac2{x-1}\right)-1,
\]
и найдите корни на $[-2;3]$.
\subhead{Ваше решение} Решение не записано. \subhead{Ваш ответ} Не указан.
\errorlabel{решение отсутствует.}
\subhead{Где именно возникает ошибка}
Не замечена замена, сводящая уравнение к квадратному.
\subhead{Почему это неверно}
При $x\ne1$ удобно ввести
\[
u=\frac{x-1}{4}-\frac2{x-1}.
\]
Тогда левая часть равна $2u^2+2$.
\keyidea{сделать замену и затем решить два рациональных уравнения.}
\subhead{Исправление и полное решение}
\[
2u^2+2=7u-1,
\]
\[
2u^2-7u+3=0,
\]
откуда $u=\frac12$ или $u=3$.

Пусть $t=x-1$. При $u=\frac12$:
\[
\frac t4-\frac2t=\frac12
\Rightarrow t^2-2t-8=0,
\]
\[
t=4\ \text{или}\ -2,
\qquad x=5\ \text{или}\ -1.
\]
При $u=3$:
\[
\frac t4-\frac2t=3
\Rightarrow t^2-12t-8=0,
\]
\[
t=6\pm2\sqrt{11},
\qquad x=7\pm2\sqrt{11}.
\]
На отрезке $[-2;3]$ лежат
\[
-1,\qquad 7-2\sqrt{11}.
\]
\subhead{Проверка}
$7-2\sqrt{11}\approx0{,}37$; остальные положительные корни больше $3$.
\finalnote{все корни: $-1$, $5$, $7\pm2\sqrt{11}$; на $[-2;3]$: $-1$ и $7-2\sqrt{11}$.}
\subhead{Задача на закрепление}
Решите
\[
\frac{(x-2)^2}{18}+\frac{18}{(x-2)^2}
=5\left(\frac{x-2}{6}-\frac3{x-2}\right)-1.
\]

\taskhead{19}
\subhead{Текст задания}
Решите
\[
\frac{3x+28-x^2}{x^2-7x}\le2+\frac2{x+3}.
\]
\subhead{Ваше решение} Решение не записано. \subhead{Ваш ответ} Не указан.
\errorlabel{решение отсутствует.}
\subhead{Где именно возникает ошибка}
Не выписана ОДЗ и не сохранена запрещённая точка $x=7$ после сокращения.
\subhead{Почему это неверно}
Исходно
\[
x\ne0,\qquad x\ne7,\qquad x\ne-3.
\]
При $x\ne7$ первая дробь сокращается, но $7$ всё равно нельзя возвращать в ответ.
\keyidea{ОДЗ записать до всех сокращений, затем применить метод интервалов.}
\subhead{Исправление и полное решение}
\[
3x+28-x^2=-(x-7)(x+4),
\]
поэтому
\[
-\frac{x+4}{x}-2-\frac2{x+3}\le0.
\]
После приведения к общему знаменателю:
\[
-\frac{3(x+1)(x+4)}{x(x+3)}\le0,
\]
или
\[
\frac{(x+1)(x+4)}{x(x+3)}\ge0.
\]
Критические точки: $-4,-3,-1,0$. Метод интервалов:
\[
(-\infty;-4]\cup(-3;-1]\cup(0;+\infty).
\]
Исключаем исходно запрещённый $x=7$:
\[
(-\infty;-4]\cup(-3;-1]\cup(0;7)\cup(7;+\infty).
\]
\subhead{Проверка}
$-4$ и $-1$ включаются как нули числителя; $-3,0,7$ исключаются.
\finalnote{ответ: $(-\infty;-4]\cup(-3;-1]\cup(0;7)\cup(7;+\infty)$.}
\subhead{Задача на закрепление}
Решите
\[
\frac{2x+15-x^2}{x^2-5x}\le1+\frac3{x+2}.
\]

\taskhead{20}
\subhead{Текст задания}
Две окружности радиусов $8$ и $15$ касаются внутренним образом в $C$. В равнобедренном прямоугольном треугольнике $ABC$ прямой угол при $C$; $A$ лежит на большей окружности, $B$ --- на меньшей. Прямая $AC$ вторично пересекает меньшую окружность в $D$, а $BC$ вторично пересекает большую в $E$. Докажите $AE\parallel BD$ и найдите $AC$.
\subhead{Ваше решение} Решение не записано. \subhead{Ваш ответ} Не указан.
\errorlabel{решение отсутствует.}
\subhead{Где именно возникает ошибка}
Не использована гомотетия с центром $C$ и не замечено, что $BD$ --- диаметр меньшей окружности.
\subhead{Почему это неверно}
При внутреннем касании гомотетия с центром $C$ и коэффициентом $\frac{15}{8}$ переводит меньшую окружность в большую:
\[
D\mapsto A,\qquad B\mapsto E.
\]
Поэтому
\[
\frac{CA}{CD}=\frac{CE}{CB}=\frac{15}{8}.
\]
\keyidea{сначала получить отношения длин из гомотетии, затем использовать прямоугольные треугольники.}
\subhead{Исправление и полное решение}
Из
\[
\frac{CB}{CD}=\frac{CE}{CA}=\frac{15}{8}
\]
и прямоугольности треугольников $BCD$ и $ECA$ имеем
\[
\tan\angle BDC=\frac{BC}{CD}
=\frac{CE}{CA}
=\tan\angle EAC.
\]
Оба угла острые, значит $\angle BDC=\angle EAC$, а так как $DC$ и $AC$ лежат на одной прямой,
\[
BD\parallel AE.
\]

Пусть $AC=BC=s$. Тогда
\[
CD=\frac8{15}s.
\]
Поскольку $\angle DCB=90^\circ$ и $D,C,B$ лежат на меньшей окружности, $DB$ --- её диаметр:
\[
DB=16.
\]
По теореме Пифагора:
\[
\left(\frac{8s}{15}\right)^2+s^2=16^2,
\]
\[
\frac{289s^2}{225}=256,
\qquad s=\frac{240}{17}.
\]
\subhead{Проверка}
При $s=\frac{240}{17}$ получаем $CD=\frac{128}{17}$, и
\[
CD^2+CB^2=16^2.
\]
\finalnote{$AE\parallel BD$, $AC=\frac{240}{17}$.}
\subhead{Задача на закрепление}
Две окружности радиусов $5$ и $12$ касаются внутренним образом в $C$. В той же конфигурации докажите параллельность соответствующих отрезков и найдите $AC$.

\taskhead{21}
\subhead{Текст задания}
Найдите все $a$, при которых
\[
\sqrt{4^x-a}+\frac{a-3}{\sqrt{4^x-a}}=1
\]
имеет ровно два различных корня.
\subhead{Ваше решение} Решение не записано. \subhead{Ваш ответ} Не указан.
\errorlabel{решение отсутствует.}
\subhead{Где именно возникает ошибка}
Не введена положительная переменная $t=\sqrt{4^x-a}$.
\subhead{Почему это неверно}
Из-за знаменателя $t>0$. После замены каждому положительному $t$ соответствует ровно одно $x$.
\keyidea{исследовать квадратное уравнение на два различных положительных корня.}
\subhead{Исправление и полное решение}
\[
t+\frac{a-3}{t}=1,
\qquad t>0,
\]
\[
t^2-t+a-3=0.
\]
Из этого же равенства
\[
t^2+a=t+3=4^x,
\]
поэтому каждому $t>0$ соответствует ровно одно
\[
x=\log_4(t+3).
\]
Нужно два различных положительных корня квадратного уравнения. Условия:
\[
D=1-4(a-3)=13-4a>0
\Rightarrow a<\frac{13}{4},
\]
и произведение корней
\[
a-3>0
\Rightarrow a>3.
\]
Итак,
\[
3<a<\frac{13}{4}.
\]
\subhead{Проверка}
При $a=3$ один корень $t=0$ недопустим; при $a=\frac{13}{4}$ корни совпадают.
\finalnote{ответ: $3<a<\frac{13}{4}$.}
\subhead{Задача на закрепление}
Найдите все $a$, при которых
\[
\sqrt{9^x-a}+\frac{a-5}{\sqrt{9^x-a}}=1
\]
имеет ровно два различных корня.

\taskhead{22}
\subhead{Текст задания}
На доске записано $10$ различных натуральных чисел. Среднее любых трёх, четырёх, пяти или шести чисел --- целое. Одно число равно $30021$.
а) Может ли быть $351$?
б) Может ли отношение двух чисел равняться $11$?
в) Если отношение двух чисел --- целое $n$, найдите минимально возможное $n$.
\subhead{Ваше решение} Решение не записано. \subhead{Ваш ответ} Не указан.
\errorlabel{решение отсутствует.}
\subhead{Где именно возникает ошибка}
Не выведено, что все числа сравнимы по модулю $60$.
\subhead{Почему это неверно}
Для любого $k\in\{3,4,5,6\}$ суммы любых $k$ чисел делятся на $k$. Сравнивая два набора, отличающиеся одним элементом, получаем: разность любых двух записанных чисел делится на $k$. Следовательно, она делится на
\[
\operatorname{lcm}(3,4,5,6)=60.
\]
\keyidea{все числа имеют один остаток по модулю $60$.}
\subhead{Исправление и полное решение}
Так как
\[
30021=60\cdot500+21,
\]
все записанные числа имеют вид $60m+21$.

\textbf{а)}
\[
351=60\cdot5+51,
\]
поэтому $351$ быть не может.

\textbf{б)}
Если $q\equiv21\pmod{60}$, то
\[
11q\equiv11\cdot21=231\equiv51\pmod{60},
\]
поэтому отношение $11$ невозможно.

\textbf{в)}
Если $nq$ и $q$ оба записаны, то
\[
nq\equiv q\pmod{60}.
\]
Так как $q\equiv21\pmod{60}$,
\[
21(n-1)\equiv0\pmod{60}.
\]
Поскольку $\gcd(21,60)=3$,
\[
20\mid(n-1).
\]
Минимальное целое $n>1$:
\[
n=21.
\]
Достижимость есть: $21$ и $441$ имеют остаток $21$ по модулю $60$, и $441/21=21$; остальные числа можно выбрать различными с тем же остатком, включая $30021$.
\subhead{Проверка}
Одинаковый остаток по модулю $60$ гарантирует целочисленность средних для $3,4,5,6$ чисел.
\finalnote{а) нет; б) нет; в) $21$.}
\subhead{Задача на закрепление}
На доске записано $8$ различных натуральных чисел. Среднее любых двух, трёх или четырёх чисел --- целое. Одно из чисел равно $1001$. Найдите минимальное возможное целое отношение большего записанного числа к меньшему, если оно больше $1$.

\vfill
\begin{center}
\textit{Лёвин Артём Александрович эксклюзивно для Ярослава Гаврилова}
\end{center}

\end{document}
